% !TEX root = ../main.tex
\lecture{7}{2026-11-12}{Digitale Güter}

\section{Digitale Güter}

\subsection{Case Android iOS}
    Welche Eigenschaft von Digitalen Güter ermöglicht starke Marktanteildominanz gewisser Anbieter
    \begin{itemize}
        \item Tendenz zu öffentliche Güter (Niedrige Eintrittsbarrieren?)
        \item Parallele zu öffentliche Güter weil Konsum von Software nicht beschränkt
        \item Monopol Tendenz
        \begin{itemize}
            \item Einführungskosten (Fix kosten) hoch aber Grenzkosten für die Bereitsstellung, Distribution niedrig
            \item Bei Wettbewerb mit neuen Konkurrenten, die vom bereits gemachten Fortschritt profitieren und kleinere Fix Kosten hatten
            \item Tendenz Preis 
        \end{itemize}
        \item Asymmetrische Information?
        \begin{itemize}
            \item Nein, es besteht nur wenn es unmöglich vor dem Konsum zu wissen, was für einen Inhalt der Ware man bekommt
            \item Zum Beispiel Qualität von Filme und Zeitung 
        \end{itemize}
        \item Netzwerkeffekte
        \begin{itemize}
            \item Steigerung der Gewinnmöglichkeiten durch Verbreitung eines Digitalen Gutes
            \item Je mehr menschen sich in einer Plattform befinden, desto mehr lohnt es sich für andere dort zu sein
            \item Voraussetzungen
            \begin{enumerate}
                \item Gebrauchswert von Produkt/Plattform abhängig der Anzahl Nutzer
                \item Aufwand als Kunde, um auf verschiedene Plattformen zu sein muss gross sein
                \item Wechselkosten von einer Plattform auf einer anderen zu migrieren müssen hoch für Konsumenten sein
            \end{enumerate}
        \end{itemize}
    \end{itemize}

\subsubsection{Netzwerkeffekte Winner-takes-all}

\begin{itemize}
    \item First mover advantage
    \begin{itemize}
        \item Erster Anbieter fängt zuerst an Kundschaftsbasis zu bilden und seinen Netzwerkeffekt erhöhen
    \end{itemize}
    \item Starke Netzwerkeffekte
    \item Multi-Homing Kosten
    \item Limitierte Produktdifferenzierung
    \item Natürliches Monopol 
    \begin{itemize}
        \item Hohe Markteintrittskosten 
    \end{itemize}
\end{itemize}

\subsection{Marktverzerrungen}
\subsubsection{ Kosten asymmetrischer Information}
Beispiel Lemon-Markets. Peach = Qualitätiv überlegene Produkte, Lemons = Gegeteil Peach.
Wenn es nicht sicher ist ob sich qualitativ besseres oder schlechteres Produkt sich vor dem Kunden befindet,
 geht man vom Erwartungswert der bekannten Verkaufspreise aus. Der Erwartungswert bei dieser Unsicherheit ist aber in manchen Fällen immer niedriger als der
 Preis des qualitativ Besseren Produkts und es werden nur Lemons gekauft. Markt für Peaches geht unter wegen Unsicherheit.
 
\begin{itemize}
    \item Versuch die Unsicherheit zu mindern
    \begin{itemize}
        \item Reviews (teilweise vom Hersteller/Anbieter gezahlt)
        \item Garantien
        \item Beschreibungen
    \end{itemize}
\end{itemize}

\subsubsection{Gegenmassnahmen}

\begin{itemize}
    \item Klagen
    \item Regulierungen
    \begin{itemize}
        \item Bei Glasfasernetz dürfen die letzte Mile des Netzes eines anderen Anbieters verwenden
    \end{itemize}
    \item Standard-Setzungen
    \begin{itemize}
        \item 
    \end{itemize}
\end{itemize}

\subsubsection{Take away}
Informationsprodukte kosten viel in der Herstellung, Profite werden von nur einigen realisiert.